\section{Reproducibility and relation to prior work}

\subsection{Provenance and workflow packaging}

W3C PROV supplies a domain-independent model of entities, activities and agents
\citep{moreau2013prov}. CWLProv maps workflow executions into interoperable
provenance records \citep{khan2019cwlprov}; RO-Crate packages research artifacts
and contextual metadata \citep{soilandreyes2022rocrate}; and Workflow Run RO-Crate
extends that packaging to multiple granularities of execution provenance
\citep{leo2024workflowrun}. This paper does not replace those models. It isolates a
narrower boundary: which additional trust, terminal and independence predicates
are required before execution records support specific research claims.

\subsection{Attestation and supply-chain integrity}

in-toto protects a software supply chain by linking authorized steps and materials
\citep{torresarias2019intoto}. SLSA defines provenance for artifacts produced by a
build platform \citep{slsa2026provenance}, and Sigstore combines identity-based
signing with transparency infrastructure \citep{sigstore2026signing}. These systems
motivate authenticated, scoped receipts and substitution attacks. The distinction
is that a scientifically admissible result needs more than an authentic execution:
it also needs a frozen evidence and adjudication protocol.

\subsection{Computational reproducibility and prospective control}

ReproZip captures dependencies and environments for computational experiments
\citep{chirigati2016reprozip}; broader reproducibility guidance calls for available
and cited data, code and workflows \citep{stodden2016enhancing}; and ACM artifact
badging separates artifact properties from reproduced or replicated results
\citep{acm2020badging}. Preregistration and Registered Reports address analytic and
publication decisions made after outcomes are known
\citep{nosek2018preregistration,chambers2015registered}. The present framework
connects these lines through a fail-closed rule: the receipt may support only the
claim whose protocol, identities and independence predicates it actually binds.

\subsection{Artifact requirements for this paper}

A submission package must contain complete tracked TeX chapters, the bibliography,
the canonical build command and dependency versions, a source-closure manifest,
the protocol and claim-ledger digests, the attack fixtures after protected release,
the raw receipt set, verification and adjudication code, and a source-to-PDF text
equivalence report. A clean build must not depend on an untracked cache or runtime
network access.

The current manuscript is a framework version. It contains no protected receipt
set and therefore makes no implementation-performance or reliability claim. This
status must remain synchronized across the abstract, Results chapter, claim ledger,
README and generated PDF.

\section{Conclusion}

Research-execution records are valuable only when their authority is stated
precisely. This paper separates request identity, result integrity, executor
attribution, evidence admission and scientific authority; defines a six-level
independence taxonomy; and specifies hostile tests that can falsify each guarantee.
The formal contract prevents typed capability failures from entering scientific
evidence and prevents an agreement label from exceeding its weakest verified gate.
These are framework properties, not a reliability result. The immediate empirical
task is therefore not to claim successful independent research execution, but to
freeze and run the adversarial protocol under authenticated, separated custody.

